\begin{theorem}[Wild quadratic phase preparation]
\label{mod:cases-wildquadratic-phasepreparation}
Construct the high- or low-range q-free common coefficient view at the actual
conductor, together with the computational data, actual phase data, exact
quadratic assembly, four actual stationary rows, and common coordinate
compatibility.  More precisely, the public type
\texttt{WildQuadraticPhasePrepared} classifies the actual coefficient input
by the proved exhaustive, disjoint dichotomy
\[
 \text{one of the three all-even rows}
 \quad\text{or}\quad
 \text{an actual \texttt{PositivePolarSource}}.
\]
Its \texttt{allEven} constructor contains the all-even coefficient view,
actual-even row compatibility, common higher data, and
\texttt{CompleteCorrectionData}; these feed
\texttt{WildQuadraticHigherData.allEven} directly.  Its
\texttt{needsRefinement} constructor contains the common view, the very
positive-polar source selected by the actual quotient-derived row, the common
higher data, the same four actual rows, and a finished q-polymorphic
\texttt{WildQuadraticPositivePolarContinuation}.  The continuation retains
the upper coordinate transport, the distinct lower-product certificate, and
all four actual phase-source equalities.  This node does not choose a
normalized refinement.

High and low are alternative internal constructions of one q-free phase
package.  The strict-high branch uses the live high-product-backed exact
assembly, while the low branch uses the live table-backed exact assembly.
Both use the actual selected \(\chi_F\) and \(\psi_F\) data and preserve every
stationary representative, row orientation, affine translation, and
\(A,u,n,z_0,z_1\) identity.
\LeanFile{LanglandsFirstMainLemma/Cases/WildQuadratic/PhasePreparation.lean}
\lean{LanglandsFirstMainLemma.wildQuadraticPhasePreparation}
\leanok
\SourceText{Propositions prop:quadratic-critical-functions, prop:quadratic-gamma-table, and prop:quadratic-complete-assembly.}
\uses{mod:ramification-primecyclicpreparation,mod:parameters-phasereduction,mod:parameters-highquadratic,mod:parameters-low,mod:cases-wildquadratic-coefficienttable,mod:cases-wildquadratic-errorformula,mod:cases-wildquadratic-coordinatetransport,mod:cases-wildquadratic-lowcoordinatetransport,mod:cases-wildquadratic-lowproductcoordinatetransport,mod:cases-wildquadratic-highcoordinatetransport,mod:parameters-minimalorbitstationary,mod:cases-wildquadratic-main}
\FormalizationContract
Produce common coefficient, phase, actual-row, and coordinate data before
any refinement is chosen.  The computational datum uses the selected base
additive character, norm pullback, trace pullback, norm-character data, and
actual minimal-orbit twists.  Its four phase rows occur in the literal
extension--norm numerator and base--twist denominator orientation.  The
common compatibility identifies the selected base and nontrivial
norm-character data and, in the stated orientation, every
\(A,u,n,z_0,z_1\) coordinate.  The positive-polar continuation is built while
the low residual polynomial or original high representatives are still live;
later refinement therefore only selects \(q\), transports the retained
certificates, and compares the retained actual phases.

The all-even branch is exactly
\texttt{belowMEvenTEven}, \texttt{boundaryEven}, or
\texttt{aboveMEvenTEven}; all four supplied Lamprecht rows are their actual
\texttt{even} constructors, and the branch returns its q-free higher datum
directly.  The other branch retains a \texttt{PositivePolarSource} built from
an actual quotient-derived lower Hasse function; a numerical parity or phase
equality is not a substitute for that source.  No characteristic-two
refinement, refined critical source, or local refinement correction is
selected here; the retained continuation is q-polymorphic.
\end{theorem}
